\section{Kodgenomgång}
\label{sec:kod}

\subsection{Arkitektur}

Koden är byggd i lager, och pilarna i figur~\ref{fig:arkitektur} går bara åt ett håll.
Datalagret vet ingenting om modellerna, modellerna vet ingenting om diagram, och allt möts i
ett gemensamt kontrakt: datatyperna i \kod{core/types.py}. Det gör att varje del kan testas för
sig och bytas ut.

\begin{figure}[ht]
\centering
\begin{tikzpicture}[font=\footnotesize\sffamily,
  srcbox/.style={draw=black!45,rounded corners=2pt,fill=black!4,align=center,inner sep=4pt,
              minimum height=9mm,minimum width=4.5cm},
  layerbox/.style={draw=accent,rounded corners=2pt,fill=accent!5,align=center,inner sep=5pt,
              minimum width=15cm,text width=14.4cm},
  modelbox/.style={draw=accent,rounded corners=2pt,fill=white,align=center,inner sep=3pt,
              minimum width=2.75cm,minimum height=11mm},
  outbox/.style={draw=boxgreen,rounded corners=2pt,fill=boxgreen!6,align=center,inner sep=3pt,
              minimum width=3.55cm,minimum height=12mm},
  flow/.style={-{Stealth[length=5pt]},thick,draw=black!60},
  lbl/.style={font=\scriptsize\sffamily,text=black!70,inner sep=1pt}]
  % data layer and its sources
  \node[layerbox] (data) {\textbf{data/}\quad \texttt{yahoo.py},
    \texttt{riksbank.py} $\rightarrow$ \texttt{store.py} (cache och fixtures) $\rightarrow$
    \texttt{mapping.py}, \texttt{fx.py}, \texttt{universe.py} $\rightarrow$ \texttt{loader.py},
    \texttt{market\_params.py}};
  \node[srcbox,above=7mm of data] (riks) {Riksbanken\\ (riskfri ränta, valutakurser)};
  \node[srcbox,anchor=south west] (yahoo) at ($(data.north west)+(0,7mm)$)
    {Yahoo Finance\\ (bokslut, kurser, universum)};
  \node[srcbox,anchor=south east] (damo) at ($(data.north east)+(0,7mm)$)
    {Referensfil, Damodaran\\ (riskpremie, skattesats)};
  \node[layerbox,below=6mm of data,fill=white] (types) {\textbf{core/types.py}\quad
    \texttt{CompanyData}\enspace \texttt{PeerSnapshot}\enspace \texttt{MarketParameters}\enspace
    \texttt{DCFAssumptions}\enspace \texttt{ValuationResult}};
  % models inside the pipeline
  \node[modelbox,anchor=north west] (comps) at ($(types.south west)+(2.5mm,-14mm)$)
    {\textbf{comps.py}\\ jämförelsebolag};
  \node[modelbox,right=12.5mm of comps] (dcf) {\textbf{dcf.py}\\ WACC och DCF};
  \node[modelbox,right=12.5mm of dcf] (lbo) {\textbf{lbo.py}\\ LBO};
  \node[modelbox,right=12.5mm of lbo] (mc) {\textbf{monte\_carlo.py}\\ simulering};
  \coordinate (padtop) at ($(comps.north west)+(0,1.5mm)$);
  \coordinate (padbottom) at ($(mc.south east)+(0,-6mm)$);
  \node[draw=accent,dashed,rounded corners=3pt,inner xsep=2.5mm,inner ysep=2mm,
        fit=(padtop)(padbottom)(comps)(mc)] (pipe) {};
  \node[font=\footnotesize\sffamily\bfseries,text=accent,fill=white,inner sep=2pt,anchor=west]
        at ($(pipe.north west)+(3mm,0)$) {pipeline.py kör modellerna i models/};
  \node[layerbox,below=7mm of pipe,fill=white,minimum width=6cm,text width=5.6cm] (report)
    {\textbf{ValuationReport}\quad alla resultat för ett bolag};
  % outputs
  \node[outbox,anchor=north west] (viz) at ($(report.south -| types.west)+(0,-9mm)$)
    {\textbf{viz/}\\ 16 diagram, football field};
  \node[outbox,right=2.67mm of viz] (xl) {\textbf{export/excel.py}\\ arbetsbok, sju flikar};
  \node[outbox,right=2.67mm of xl] (app) {\textbf{app/streamlit\_app.py}\\ webbapp};
  \node[outbox,right=2.67mm of app] (cli) {\textbf{\_\_main\_\_.py}, \textbf{batch.py}\\
    kommandorad, alla bolag};
  % edges
  \draw[flow] (yahoo.south) to (yahoo.south |- data.north);
  \draw[flow] (riks.south) to (data.north);
  \draw[flow] (damo.south) to (damo.south |- data.north);
  \draw[flow] (data.south) to (types.south |- data.south) to (types.north);
  \draw[flow] (types.south) to (types.south |- pipe.north);
  \draw[flow] (comps.east) to node[lbl,above=1pt] {peers} (dcf.west);
  \draw[flow] (dcf.east) to node[lbl,above=1pt] {prognos} (lbo.west);
  \draw[flow] (dcf.south) to[out=-18,in=-162,looseness=0.45]
    node[lbl,below=0.5pt,pos=0.5] {basfall} (mc.south);
  \draw[flow] (pipe.south) to (report.north);
  \coordinate (bus) at ($(report.south)+(0,-4.5mm)$);
  \draw[thick,draw=black!60] (report.south) to (bus);
  \draw[flow] (bus) -| (viz.north);
  \draw[flow] (bus) -| (xl.north);
  \draw[flow] (bus) -| (app.north);
  \draw[flow] (bus) -| (cli.north);
\end{tikzpicture}
\caption{Arkitekturen: data, modeller, pipeline och presentation. Pipelinen kör comps först och
skickar jämförelsebolagen till DCF:en. LBO:n får DCF:ens prognos och exit-multipel, och
simuleringen utgår från DCF:ens basfall.}
\label{fig:arkitektur}
\end{figure}

\subsection{Modul för modul}

\begin{description}
  \item[\kod{core/types.py}] Kontraktet. Alla belopp är absoluta tal i bolagets handelsvaluta,
    kostnader lagras som positiva tal, och räntor som decimaltal. En post som saknas är
    \kod{NaN}, aldrig en saknad kolumn. Tabell~\ref{tab:typer} listar de viktigaste typerna.
  \item[\kod{data/}] \kod{yahoo.py} är den enda modul som pratar med nätet. Den sparar råa svar
    som JSON i en lokal cache, och alla objekt byggs sedan från de filerna. \kod{mapping.py}
    översätter Yahoos radnamn till standardposter genom ordnade reservkedjor, och vilken källa
    som användes sparas i \kod{source_map}. \kod{universe.py} rensar Yahoos lista från 984
    symboler till 687 bolag (ett aktieslag per bolag). \kod{fx.py} räknar om bolag som
    rapporterar i annan valuta med Riksbankens kurser. Detaljerna står i
    \kod{docs/data_mapping.md}.
  \item[\kod{models/dcf.py}] \kod{default_assumptions} bygger antagandena ur historiken,
    \kod{compute_wacc} räknar WACC med alla delsteg, och \kod{dcf_core} gör själva
    aritmetiken. \kod{dcf_core} är vektoriserad och den enda platsen där DCF-formlerna finns:
    basfallet, känslighetstabellerna och simuleringen anropar samma funktion.
  \item[\kod{models/comps.py}] \kod{run_comps} väljer jämförelsebolag enligt
    avsnitt~\ref{sec:peers-arligt}. Varje avsteg från huvudvägen loggas i
    \kod{CompsResult.fallbacks}.
  \item[\kod{models/lbo.py}] \kod{run_lbo} bygger affären, skuldschemat och avkastningen.
    \kod{compute_irr} och lösaren för betalningsförmåga ligger här.
  \item[\kod{models/monte_carlo.py}] \kod{run_monte_carlo} tar DCF:ens exakta indata
    (\kod{DCFResult.core_inputs}), byter fyra av dem mot vektorer och anropar \kod{dcf_core}.
  \item[\kod{pipeline.py}] \kod{run_valuation(ticker)} kör allt för ett bolag och returnerar
    en \kod{ValuationReport}. En modell som inte går att tillämpa kastar
    \kod{NotApplicableError} med en förklaring, och de andra modellerna körs ändå. Banker och
    fastighetsbolag får till exempel bara P/E.
  \item[\kod{batch.py}, \kod{__main__.py}] Kommandoraden. \kod{run_universe} värderar alla
    bolag och skriver en CSV-fil med en rad per bolag.
  \item[\kod{viz/}, \kod{export/excel.py}, \kod{app/}] Diagram med matplotlib i gemensam stil,
    en Excel-arbetsbok med flikarna Summary, Assumptions, DCF, Comps, LBO, Monte Carlo och
    Data (blå text för indata, svart för beräknat), och en Streamlit-app där antagandena kan
    ändras i sidopanelen.
\end{description}

\begin{table}[ht]
\centering\small
\caption{De viktigaste datatyperna.}
\label{tab:typer}
\begin{tabularx}{\linewidth}{@{}l>{\raggedright\arraybackslash}X@{}}
\toprule
Typ & Innehåll \\
\midrule
\kod{CompanyData} & Allt om ett bolag: \kod{profile} (namn, sektor, valuta), \kod{statements}
  (tre rapporter som tabeller per räkenskapsår), \kod{market} (kurs, börsvärde, skuld, kassa)
  och \kod{price_history} (veckokurser för beta) \\
\kod{PeerSnapshot} & En rad per bolag i universumet, i kronor: börsvärde, rörelsevärde,
  omsättning, EBITDA, vinst, beta. Används av comps och för beta \\
\kod{MarketParameters} & Riskfri ränta, riskpremie, skattesats, index och evig tillväxt \\
\kod{DCFAssumptions} & Tillväxt och marginal per år, avskrivningar, capex och rörelsekapital
  som andel av omsättningen, skatt, WACC, evig tillväxt, exit-multipel \\
\kod{ValuationResult} & En stapel i football field: \kod{method}, \kod{low}, \kod{base},
  \kod{high}. Varje modell har \kod{to_valuation_results()} \\
\kod{ValuationReport} & Pipelinens resultat: de fyra modellresultaten, staplarna, varningar och
  parametrarnas källor \\
\bottomrule
\end{tabularx}
\end{table}

\subsection{Installera och köra}

\begin{lstlisting}[style=sh,caption={Installation och de vanligaste kommandona.},label=lst:kor]
# installation (Python 3.11 eller senare)
python -m venv .venv && source .venv/bin/activate
pip install -e ".[dev]"

# ett bolag: skriver ut en sammanfattning och sparar diagram, Excel och summary.json
python -m valuation_engine ATCO-A.ST --offline

# hela universumet: ladda först ned data (ungefär 18 minuter), kör sedan alla bolag
python -m valuation_engine.data.download --all
python -m valuation_engine --all

# appen och testerna
streamlit run app/streamlit_app.py
pytest
\end{lstlisting}

Resultaten hamnar i \kod{output/<TICKER>/}. Användbara flaggor: \flag{simulations} och
\flag{seed} för simuleringen, \flag{no-size-premium} för ren CAPM, och \flag{universe-peers}
för att låta DCF:en hämta beta och exit-multipel från hela branschen i stället för från comps.
Samma sak från Python:

\begin{lstlisting}[caption={Motorn som bibliotek: basfallet, och caset med två ändrade
antaganden.},label=lst:api]
from dataclasses import replace
from valuation_engine.pipeline import run_valuation
from valuation_engine.models import dcf

report = run_valuation("ATCO-A.ST", offline=True)
print(round(report.dcf.value_per_share_gordon, 2))        # 115.85
for v in report.valuations:
    print(v.label, round(v.low, 1), round(v.base, 1), round(v.high, 1))

base = dcf.default_assumptions(report.company, report.params)
mine = replace(base, tax_rate=0.219, nwc_pct_revenue=0.133)
adjusted = run_valuation("ATCO-A.ST", offline=True, dcf_assumptions=mine)
print(round(adjusted.dcf.value_per_share_gordon, 1))      # 115.9
\end{lstlisting}

\subsection{Offline-läge och fixtures}

Med flaggan \flag{offline} eller miljövariabeln \kod{VALUATION_ENGINE_OFFLINE=1} rör motorn
aldrig nätet. Den läser då i första hand de frysta testdata som ligger i repot
(\kod{tests/fixtures/}) och i andra hand den lokala cachen. Fixtures innehåller femton bolag
(Atlas Copco och fjorton andra storbolag, bland dem banken SEB), hela universumlistan, tabellen
med alla jämförelsebolag och Riksbankens serier. Det räcker för att köra caset exakt, eftersom
comps bara behöver tabellen. Hela universumet kräver nedladdning, eftersom cachen inte ligger
i git. Och eftersom alla objekt byggs från sparade råfiler ger en körning offline exakt samma
tal som den körning som skapade filerna.

\subsection{Testerna}

Testerna ligger i \kod{tests/}, en fil per del: \kod{test_types}, \kod{test_data},
\kod{test_dcf}, \kod{test_wacc}, \kod{test_comps}, \kod{test_lbo}, \kod{test_monte_carlo},
\kod{test_pipeline}, \kod{test_cli_batch}, \kod{test_export_viz} och \kod{test_app}. Sviten har
omkring 650 tester och går på ungefär en minut. Tre idéer är värda att ta med sig:

\begin{itemize}
  \item \textbf{Inget nät.} \kod{conftest.py} slår på offline-läget, pekar cachen mot en tom
    katalog och får varje försök att öppna en nätverksanslutning att fälla testet. Tester som
    behöver nätet är märkta \kod{network} och körs inte som standard.
  \item \textbf{Invarianter.} \kod{invariants.py} räknar om samband som måste gälla för
    \emph{varje} bolag, ur de tabeller modellerna rapporterar: att FCFF är summan av sina delar,
    att bryggan går ihop, att källor är lika med användning, att skulden aldrig blir negativ.
    \kod{test_fixture_universe.py} kör alla modeller på alla fixtures och på en svärm av
    syntetiska bolag. En modell ska antingen ge ett resultat som klarar invarianterna eller
    kasta \kod{NotApplicableError}.
  \item \textbf{Kända facit.} Där det finns en sluten lösning testas den: noll varians i
    simuleringen ska ge basfallet exakt, och ett fast slumpfrö ska ge identiska resultat.
\end{itemize}
